\providecommand{\Cont}{\operatorname{Cont}}
\providecommand{\Symp}{\operatorname{Symp}}
\providecommand{\Ham}{\operatorname{Ham}}
\providecommand{\Leg}{\operatorname{Leg}}
\let\C\relax
\newcommand{\C}{\mathbb{C}}

\chapter{Yakov Eliashberg (2020)}
\index{Eliashberg, Yakov}

\begin{quote}
\it ``Yakov Eliashberg has been the leading figure in symplectic and contact topology for decades. His discovery of overtwisted contact structures, his development of the h-principle, and his creation of symplectic field theory have transformed our understanding of odd-dimensional geometry and its connections to analysis and physics.'' --- Wolf Prize Citation, 2020
\end{quote}

\section{Early Life and Career}

Yakov Markovich Eliashberg was born in 1946 in Leningrad, USSR (now Saint Petersburg). He studied at Leningrad State University, receiving his Ph.D.\ in 1972 under the supervision of Rokhlin. He emigrated to the United States in 1987 and has been a professor at Stanford University since 1989. He was awarded the Wolf Prize in Mathematics in 2020.

\section{Main Contributions}

Yakov Eliashberg (born 1946) is a Russian--American mathematician who has fundamentally shaped symplectic and contact geometry and topology. He was awarded the Wolf Prize in Mathematics in 2020 alongside Simon Donaldson.

The Wolf Prize citation reads: ``for their contributions to differential geometry and topology.''

Eliashberg's core contributions include:
\begin{itemize}
    \item \textbf{Overtwisted Contact Structures:} The classification of overtwisted contact structures on 3-manifolds, showing they are completely determined by algebraic topological data (the h-principle holds), in stark contrast to tight contact structures.
    \item \textbf{Flexible and Rigid Symplectic Geometry:} The discovery of the dichotomy between flexibility (h-principle) and rigidity (Gromov's theory of pseudoholomorphic curves) in symplectic topology.
    \item \textbf{Symplectic Field Theory (SFT):} The creation (with Givental and Hofer) of a comprehensive framework for invariants of symplectic and contact manifolds using punctured pseudoholomorphic curves.
    \item \textbf{Classification of Legendrian Knots:} The classification of Legendrian and transverse knots in contact 3-manifolds using invariants from SFT.
    \item \textbf{Stein Manifolds:} The characterization of Stein manifolds (complex manifolds admitting a proper plurisubharmonic function) in terms of handlebody decompositions, leading to the construction of exotic Stein structures.
    \item \textbf{The h-Principle:} Systematic development of the h-principle (homotopy principle) for differential relations, showing that many geometric problems reduce to algebraic topology.
\end{itemize}

\section{Origin of the Problem}

\subsection{Contact Structures and Flexibility}

A contact structure on an odd-dimensional manifold is a maximally non-integrable hyperplane field. It is the odd-dimensional analogue of a symplectic structure. The fundamental question: when do two contact structures on the same manifold differ by a diffeomorphism?

The discovery: there is a fundamental dichotomy. Some contact structures (``overtwisted'') are flexible and classified by homotopy theory. Others (``tight'') are rigid and require pseudoholomorphic curves for their study.

\subsection{The h-Principle}

In differential geometry, many problems involve finding a solution to a differential equation or relation. The h-principle states that the existence of a formal solution (satisfying the algebraic topological constraints) often implies the existence of an actual solution.

For example, the h-principle for immersions says that any map $f: M \to N$ with $df$ injective (formally an immersion) can be deformed to an actual immersion.

\subsection{Symplectic Field Theory}

Symplectic field theory provides a unified framework for counting pseudoholomorphic curves in symplectic and contact manifolds. It generalizes Gromov--Witten theory and Floer homology, and produces invariants of symplectic and contact structures.

\section{How It Was Solved and the Intuitions/Tools Needed}

\subsection{The Classification of Overtwisted Contact Structures}

Eliashberg proved that on a 3-manifold, every homotopy class of plane fields that are formally contact (i.e., a reduction of the structure group to $\U(1) \times 1$) contains a unique overtwisted contact structure up to isotopy. The proof uses a combination of:
\begin{enumerate}
    \item The parametric h-principle for overtwisted contact structures.
    \item A surgery decomposition of 3-manifolds.
    \item Explicit model contact structures on standard pieces.
\end{enumerate}

\subsection{Tools and Techniques}

\begin{itemize}
    \item \textbf{Pseudoholomorphic Curves\index{holomorphic curve}:} Gromov's theory of $J$-holomorphic curves is the main source of rigidity in symplectic topology.
    \item \textbf{Holomorphic Disk Filling:} A technique to detect tightness using families of pseudoholomorphic disks.
    \item \textbf{Convex Surface Theory:} The study of contact structures via convex surfaces (Giroux).
    \item \textbf{Hurwitz--Morse Theory:} The topology of Stein manifolds via Morse theory of plurisubharmonic functions.
\end{itemize}

\section{Mathematical Theory}

\subsection{Contact Structures}

\begin{definition}[Contact Structure]
A \textbf{contact structure} on a $(2n+1)$-dimensional manifold $M$ is a hyperplane field $\xi \subset TM$ that is maximally non-integrable: if $\xi = \ker \alpha$ locally, then $\alpha \wedge (d\alpha)^n \neq 0$.
\end{definition}

\begin{definition}[Overtwisted vs Tight]
A contact 3-manifold $(M,\xi)$ is \textbf{overtwisted} if it contains an overtwisted disk: an embedded disk $D$ such that $\xi$ is tangent to $\partial D$ and transverse to the interior of $D$. Otherwise, $\xi$ is \textbf{tight}.
\end{definition}

\begin{theorem}[Eliashberg, 1989]
On any closed oriented 3-manifold $M$, there is a natural bijection between:
\[
\frac{\{\text{overtwisted contact structures}\}}{\text{isotopy}} \cong \frac{\{\text{formal contact structures}\}}{\text{homotopy}}
\]
In other words, overtwisted contact structures satisfy the h-principle and are classified by algebraic topology.
\end{theorem}

\subsection{The h-Principle}

\begin{definition}[Differential Relation]
A \textbf{differential relation} $\mathcal{R}$ of order $k$ is a subset of the jet bundle $J^k(M,N)$. A solution is a function $f: M \to N$ such that $j^k f(M) \subset \mathcal{R}$.
\end{definition}

\begin{definition}[h-Principle]
We say that the \textbf{h-principle} holds for a differential relation $\mathcal{R}$ if every formal solution (a section $F: M \to \mathcal{R}$) is homotopic through formal solutions to a holonomic solution (one of the form $j^k f$ for some $f$).
\end{definition}

Eliashberg proved the h-principle for many geometric problems, including:
\begin{itemize}
    \item Immersions and submersions.
    \item Symplectic and contact embeddings.
    \item Over twised contact structures.
    \item Legendrian immersions.
\end{itemize}

\subsection{Symplectic Field Theory}

\begin{definition}[Symplectic Field Theory]
SFT is a collection of invariants of symplectic and contact manifolds defined by counting pseudoholomorphic curves in the symplectization $\R \times M$ of a contact manifold $M$, with punctures asymptotic to periodic Reeb orbits.
\end{definition}

The SFT Hamiltonian $H$ is a formal power series in variables corresponding to Reeb orbits, satisfying:
\[
H \circ H = 0
\]
making it a differential graded algebra. The homology of this complex is an invariant of the contact structure.

\subsection{Stein Manifolds}

\begin{definition}[Stein Manifold]
A complex manifold $W$ is \textbf{Stein} if it admits a proper holomorphic embedding into $\C^N$ for some $N$. Equivalently, $W$ admits a smooth proper plurisubharmonic function $\varphi: W \to \R$ with finitely many critical points.
\end{definition}

\begin{theorem}[Eliashberg, 1990]
A smooth almost complex manifold $W$ of dimension $2n$ (with $n > 2$) admits a Stein structure if and only if it has a handlebody decomposition with handles of index at most $n$.
\end{theorem}

This theorem enables the construction of exotic Stein manifolds and establishes a deep connection between complex geometry and contact topology (the boundary of a Stein manifold inherits a tight contact structure).

\subsection{Gromov's h-Principle for Contact Structures}

\begin{theorem}[Gromov, 1973; Eliashberg, 1989]
Let $M$ be a closed oriented manifold of dimension $2n+1 \geq 5$. Then overtwisted contact structures on $M$ satisfy the \textbf{h-principle}: two overtwisted contact structures are isotopic if and only if they are homotopic as plane fields. In dimension 3, Eliashberg (1989) proved that every homotopy class of formally almost contact structures contains a unique overtwisted contact structure up to isotopy.
\end{theorem}

\begin{remark}
The h-principle for overtwisted contact structures is remarkable because contact structures are defined by a nonlinear PDE (maximal non-integrability of a distribution), yet the classification reduces to a purely algebraic-topological problem. The key ingredient in Eliashberg's proof is the \textbf{parametric h-principle}: one shows that the inclusion of isotopy classes of overtwisted contact structures into homotopy classes of almost contact structures is a bijection by constructing explicit isotopies using surgery techniques.
\end{remark}

\begin{remark}
In dimensions $\geq 5$, the result was proved by Gromov using his general theory of convex integration. The dimension 3 case required fundamentally new ideas, as the topology of 3-manifolds allows for local modifications (Legendrian surgery) that are not available in higher dimensions.
\end{remark}

\subsection{Fillability and the Weinstein Conjecture}

\begin{definition}[Symplectic Fillability]
A contact manifold $(M, \xi)$ is \textbf{symplectically fillable} if there exists a symplectic manifold $(W, \omega)$ with $\partial W = M$ (as oriented manifolds) such that $\omega|_{\xi}$ is a positive volume form on $\xi$ at each point of $\partial W$. It is \textbf{exact fillable} if additionally $\omega = d\lambda$ with $\lambda|_M$ a contact form for $\xi$.
\end{definition}

\begin{definition}[Weinstein Conjecture]
Let $(M, \xi)$ be a closed contact manifold with a contact form $\alpha$ (so $\xi = \ker \alpha$). The \textbf{Reeb vector field} $R_\alpha$ is defined by $\iota_{R_\alpha} d\alpha = 0$ and $\alpha(R_\alpha) = 1$. The Weinstein conjecture asserts that the Reeb vector field of any contact form on a closed contact manifold has a periodic orbit.
\end{definition}

\begin{theorem}[Hofer, 1993]
Every closed overtwisted contact 3-manifold satisfies the Weinstein conjecture: the Reeb vector field of any contact form has a periodic orbit.
\end{theorem}

\begin{remark}
The relationship between fillability and the Weinstein conjecture is deep. For overtwisted contact manifolds, the h-principle implies fillability: every overtwisted contact structure on a closed 3-manifold is fillable (by a theorem of Eliashberg). However, there exist tight contact structures that are not fillable, providing a sharp distinction between the flexible and rigid categories. The Weinstein conjecture is now known to hold for all contact 3-manifolds by the work of Colin--Honda (2008), which builds on Eliashberg's foundational classification.
\end{remark}

\begin{remark}[Exact Lagrangian Submanifolds]
A Lagrangian submanifold $L \subset (W, \omega)$ is \textbf{exact} if $\omega|_L = d\beta$ for a 1-form $\beta$ on $L$ that extends to a primitive $\lambda$ with $\omega = d\lambda$ on $W$. Exact Lagrangian submanifolds cannot be closed (by Stokes' theorem), but their existence imposes strong constraints on fillability. Gromov (1985) showed that a closed exact Lagrangian submanifold cannot exist in an exact symplectic manifold. This result connects the topology of Lagrangian submanifolds to the fillability of contact boundaries via the exact sequence relating Legendrian knots in $\partial W$ to Lagrangian submanifolds in $W$.
\end{remark}

\subsection{Flexible and Rigid Contact Structures}

\begin{definition}[Flexible (Overtwisted) Contact Structure]
A contact structure $\xi$ on a manifold $M$ is \textbf{flexible} (or \textbf{overtwisted} in dimension 3) if it satisfies the h-principle: its isotopy class is determined by its homotopy class as a plane field. Flexible structures admit arbitrarily large $C^0$-perturbations and are classified by algebraic topology. In dimension 3, this is equivalent to being overtwisted.
\end{definition}

\begin{definition}[Rigid (Tight) Contact Structure]
A contact structure $\xi$ is \textbf{tight} if it does not contain an overtwisted disk (in dimension 3) or, more generally, if it satisfies certain rigidity properties. Tight structures require the full apparatus of pseudoholomorphic curve theory for their study and can distinguish between contact structures that are homotopic but not isotopic.
\end{definition}

\begin{theorem}[Eliashberg's Dichotomy]
On a closed 3-manifold, every contact structure is either:
\begin{enumerate}
    \item \textbf{Overtwisted:} classified by homotopy theory (flexible), or
    \item \textbf{Tight:} not isotopic to any overtwisted structure (rigid).
\end{enumerate}
Moreover, the standard contact structure on $S^3$ (the kernel of the connection 1-form) is tight, and any contact structure isotopic to an overtwisted structure is itself overtwisted.
\end{theorem}

\begin{remark}[Characterization of Tightness via Holomorphic Curves]
The deep reason tight structures are rigid is that Gromov's theory of pseudoholomorphic curves provides powerful invariants. In a tight contact manifold, the moduli spaces of holomorphic disks with boundary on Legendrian submanifolds carry transversality properties that are absent in the overtwisted case. Specifically, tightness prevents the existence of ``fat'' Legendrian knots---those bounding a disk with too many Reeb chord intersections---which is precisely the obstruction captured by the overtwisted disk.
\end{remark}

\begin{remark}[Higher-Dimensional Flexibility and Rigidity]
In dimensions $\geq 5$, the overtwisted/tight dichotomy does not generalize in the same way. Instead, the relevant dichotomy is between:
\begin{itemize}
    \item \textbf{Flexible contact structures:} those satisfying the h-principle (all contact structures in dimension $\geq 5$ satisfying certain local conditions).
    \item \textbf{Strongly fillable (Weinstein) structures:} those admitting exact symplectic fillings, which carry the full force of holomorphic curve rigidity.
\end{itemize}
This higher-dimensional theory was developed by Borman--Eliashberg--Murphy (2015), who proved the h-principle for flexible contact structures in all dimensions.
\end{remark}

\subsection{The Eliashberg--Gromov Dichotomy in Symplectic Geometry}

\begin{definition}[Symplectic Rigidity]
A symplectic manifold $(W, \omega)$ is \textbf{rigid} if its symplectic diffeomorphism group acts discretely on the space of symplectic structures---meaning small perturbations of $\omega$ cannot be absorbed by a diffeomorphism. This is the generic behavior for closed symplectic manifolds and is detected by Gromov's pseudoholomorphic curve invariants.
\end{definition}

\begin{definition}[Symplectic Flexibility]
A symplectic structure is \textbf{flexible} if the symplectic embedding problem reduces to volume considerations. Eliashberg's characterization (2001) of flexible symplectic manifolds shows that ``weinstein Weinstein''---a symplectic manifold is flexible if and only if it is Weinstein and its Weinstein structure can be deformed to a particularly simple form.
\end{definition}

\begin{theorem}[Eliashberg, 2001]
A Stein domain $(W, J)$ of dimension $\geq 6$ is flexible if and only if its Weinstein structure can be deformed so that the core Legendrian submanifolds have no ``exotic'' holomorphic disk counts. In dimension 4, flexibility is equivalent to the vanishing of certain Gromov--Witten-type invariants.
\end{theorem}

\begin{remark}[Contact Surgery and Flexibility]
The transition from flexible to rigid structures is controlled by Legendrian surgery: cutting along a Legendrian sphere and gluing in a Weinstein handle. Each such surgery either preserves flexibility or creates rigidity, depending on the holomorphic disk count on the attaching region. This perspective unifies the classification of contact structures (flexible = classified by h-principle) with the classification of Stein fillings (rigid = classified by holomorphic curves).
\end{remark}

\begin{remark}[Symplectic Homology as an Obstruction]
Symplectic homology $\mathrm{SH}(W)$, introduced by Viterbo and refined by Eliashberg--Givental--Hofer, is a powerful invariant that vanishes for flexible (exact) Weinstein manifolds but is typically non-vanishing for rigid ones. It provides a computable obstruction to flexibility and detects the failure of the h-principle for symplectic embeddings.
\end{remark}

\section{Impact on Science and Technology}

\subsection{In Mathematics}

\begin{enumerate}
    \item \textbf{Symplectic Topology:} The flexibility/rigidity dichotomy is fundamental.
    \item \textbf{Contact Topology:} The classification of overtwisted structures is the central result in the field.
    \item \textbf{Stein Geometry:} Eliashberg's characterization enabled the construction of many new Stein manifolds.
    \item \textbf{Topology of 3-Manifolds:} Contact structures are deeply connected to the geometry of 3-manifolds.
\end{enumerate}

\subsection{In Physics}

\begin{enumerate}
    \item \textbf{Symplectic Field Theory:} SFT is related to topological field theories and string theory.
    \item \textbf{Classical Mechanics:} Contact geometry\index{contact geometry} appears in thermodynamics and geometric optics.
    \item \textbf{Quantization:} Contact structures are related to the prequantization of symplectic manifolds.
\end{enumerate}

\section{Frequently Asked Questions}

\subsection*{Q1: What is an overtwisted contact structure?}
A contact structure on a 3-manifold that contains an overtwisted disk—a disk whose boundary is tangent to the contact planes. Overtwisted structures are flexible and classified by algebraic topology.

\subsection*{Q2: What is the h-principle?}
The homotopy principle states that for certain differential relations, the existence of a formal (algebraic topological) solution implies the existence of an actual geometric solution. It is the source of flexibility in geometry.

\subsection*{Q3: What is symplectic field theory?}
A comprehensive framework for invariants of symplectic and contact manifolds, counting punctured pseudoholomorphic curves in symplectizations. It generalizes Gromov--Witten theory and Floer homology.

\subsection*{Q4: What is a Stein manifold?}
A complex manifold that can be properly holomorphically embedded in $\C^N$. Eliashberg characterized Stein manifolds as those with a handlebody decomposition where handles have index at most the complex dimension.

\subsection*{Q5: What should an undergraduate study to understand Eliashberg's work?}
Differential topology (manifolds, vector bundles, transversality), symplectic geometry (symplectic forms, Hamiltonian dynamics), and complex geometry. Recommended: \textit{Introduction to Symplectic Topology} (McDuff--Salamon), \textit{Contact Geometry} (Geiges), and \textit{The h-Principle} (Eliashberg--Mishachev).

\section{Related Chapters}
See also Chapter~22 (Milnor) on differential topology.

\section*{Exercises for the Reader}
\addcontentsline{toc}{section}{Exercises}

\begin{enumerate}
    \item [B] Show that $\R^3$ with $\xi = \ker(dz + x\,dy)$ is a contact structure.
    \item [I] Prove that an overtwisted disk cannot exist in the standard tight contact structure on $S^3$.
    \item [A] Show that every closed oriented 3-manifold admits an overtwisted contact structure.
    \item [I] Prove the h-principle for immersions of $S^1$ into $\R^2$.
    \item [I] Show that the boundary of a Stein manifold carries a natural contact structure.
    \item [A] Construct an exotic Stein structure on $\R^4$.
    \item [B] Show that the Reeb vector field of a contact form generates a volume-preserving flow.
\end{enumerate}

\section*{Timeline}
\addcontentsline{toc}{section}{Timeline}

\begin{tabularx}{\linewidth}{|p{1.5cm}|>{\raggedright\arraybackslash}X|}
\hline
\textbf{Year} & \textbf{Event} \\ \hline
1946 & Born in Leningrad, USSR \\ \hline
1968 & M.S.\ from Leningrad State University \\ \hline
1972 | Ph.D.\ from Leningrad State University \\ \hline
1988 | Moves to the United States (Stanford) \\ \hline
1989 | Classification of overtwisted contact structures \\ \hline
1990 | Characterization of Stein manifolds \\ \hline
1994 | Moves to Stanford University \\ \hline
2000 | Symplectic field theory (with Givental and Hofer) \\ \hline
2005 | Veblen Prize of the AMS \\ \hline
2015 | Moves to Stanford (continuing) \\ \hline
2020 | Wolf Prize in Mathematics \\ \hline
\end{tabularx}

\section*{Glossary}
\addcontentsline{toc}{section}{Glossary}

\begin{description}
    \item[Contact structure] A maximally non-integrable hyperplane field on an odd-dimensional manifold.
    \item[h-Principle] The principle that formal solutions imply actual geometric solutions.
    \item[Legendrian knot] A knot tangent to a contact structure.
    \item[Overtwisted] A contact structure containing an overtwisted disk.
    \item[Pseudoholomorphic curve] A map from a Riemann surface to an almost complex manifold satisfying the Cauchy--Riemann equations.
    \item[Reeb orbit] A periodic orbit of the Reeb vector field of a contact form.
    \item[Stein manifold] A complex manifold with a proper plurisubharmonic function.
    \item[Symplectic field theory] An invariant counting punctured pseudoholomorphic curves in symplectizations.
    \item[Tight] A contact structure that is not overtwisted.
\end{description}

\section{Citations}\subsection*{Primary Sources}
\begin{enumerate}
    \item Y.\ Eliashberg, \textit{Classification of overtwisted contact structures on 3-manifolds}, Invent.\ Math.\ 98 (1989), 623--637.
    \item Y.\ Eliashberg, \textit{Topological characterization of Stein manifolds of dimension $> 2$}, Int.\ J.\ Math.\ 1 (1990), 29--46.
    \item Y.\ Eliashberg, A.\ Givental, and H.\ Hofer, \textit{Introduction to symplectic field theory}, GAFA 2000 (Tel Aviv), 560--673.
    \item Y.\ Eliashberg and N.\ Mishachev, \textit{Introduction to the h-Principle}, AMS Grad.\ Stud.\ in Math.\ 48, 2002.
    \item Y.\ Eliashberg, \textit{Symplectic geometry and topology}, Proc.\ ICM 1998.
\end{enumerate}

\subsection*{Expository Works}
\begin{enumerate}
    \item H.\ Geiges, \textit{An Introduction to Contact Topology}, Cambridge Univ.\ Press, 2008.
    \item D.\ McDuff and D.\ Salamon, \textit{Introduction to Symplectic Topology}, Oxford Univ.\ Press, 2017.
    \item K.\ Cieliebak and Y.\ Eliashberg, \textit{From Stein to Weinstein and Back}, AMS Colloquium Publ., 2012.
\end{enumerate}

